We presented an evidence-constrained workflow for diagnostic assistance on unlabeled neutron-current time series. The workflow uses deterministic analysis probes under a persistent workspace contract to produce auditable observations, then turns a structured alert summary into manual-grounded, citation-bearing explanatory support. Its output is deliberately framed as an aid to engineering review rather than an autonomous fault decision.

The retrieval study identifies a practical governance issue: for the evaluated corpus and embedding model, cosine similarity can improve while the priority of pre-designated IAEA sources declines as same-domain supplementary manuals are added. Short Chinese queries are particularly sensitive in the tested setting, and document-level relevance is more forgiving than exact chunk identity under overlapping segmentation. These results favor terminology-rich report-derived queries and explicit source policies over similarity-only deployment choices.

Future work will validate the workflow with expert review of observed phenomena and citation support, construct a small relevance-annotated query set, and assess multilingual encoders, rerankers, metadata controls, and generation-side factuality. Such evaluation is required before any operational use beyond supervised diagnostic assistance.
